% !TEX root = ../manuscript.tex
\section{Conclusions}
\label{sec:conclusions}

This paper proposed SC-FMA as a representation layer that turns supplied utility fields or annotation-derived fidelity fields into inspectable audit records for knowledge-intensive artifacts. The real-data controls show that graph topology contributes little independent PRM800K annotation-order signal, while windowing can partially repair the long-trace failure of the full QP allocation.

Across complementary checks, SC-FMA turns supplied fidelity and structural fields into audit records with explicit audit reasons. The decisive real-data result is not a performance improvement: on the locked PRM800K route, Ridge closely tracks but does not improve over the supervised \wstruct{} fidelity field. A same-supervision graph-only Ridge reaches $\rho=0.043$, and direct mathematical-DAG necessity remains below a reverse-position baseline. These findings make Ridge the conservative reportable view and graph fields organizational rather than independently predictive on this distribution. Unwindowed QP degrades on long traces; a post hoc windowed diagnostic raises the long-stratum result from $\rho=0.172$ to $0.385$ but does not surpass the fidelity-tracking references. Redundancy and bottleneck terms act as conditional regularizers that contribute primarily in designed stress-test settings or when fidelity fields lack structural preference.

The claim remains bounded to observable artifacts and fixed-budget representation checks. Within that boundary, SC-FMA provides a reusable audit-record schema for systems where intermediate knowledge artifacts must be inspected, maintained, and reused under limited curation budgets. Four extensions follow directly from this boundary. First, domain-specific graph backends should be tested on settings where dependency structure is independently annotated rather than inferred from step position. Second, dynamic-budget updating should test how audit records change when review capacity shifts during operation. Third, human-in-the-loop validation should expose SC-FMA records to real domain experts and measure interpretability, actionability, timing, and disagreement with human audit choices. Fourth, cross-domain transfer should test the schema on RAG evidence traces, knowledge-graph maintenance records, software-engineering logs, and other knowledge-intensive workflows; long-trace validation should select sparse, blockwise, or windowed calibration on development data before independent reporting.

